\section{KL Measurement Basics}
\label{sec:kl-basics}

In traditional weak lensing (WL), we measure the ellipticity of galaxies $\hat{\bm{\epsilon}}\equiv\hat{\epsilon}_1+i\hat{\epsilon_2}$, which is a combination of the galaxy's intrinsic shape $\hat{\bm{\epsilon}}^\mathrm{int}$ (primarily due to inclination) and reduced shear $\bm{g}\equiv g_1+ig_2$ \citep{WL2001}:
\begin{equation}
\label{eqn:shear-transform}
    \hat{\bm{\epsilon}} = \frac{\hat{\bm{\epsilon}}^\mrm{int}+\bm{g}}{1+\bm{g}^*\hat{\bm{\epsilon}}^\mathrm{int}}.
\end{equation}
Since $\hat{\bm{\epsilon}}$ is correlated with $\bm{g}$, it is often used as a proxy estimator for $\bm{g}$. Along with the fact that intrinsic shapes are uncorrelated over large distances, and that shear effects are typically very small, we can approximate Eq.~(\ref{eqn:shear-transform}) as a linear relationship to calibrate the estimator on an ensemble measurement:
\begin{equation}
\label{eqn:linear-bias}
    \langle\hat{\bm{\epsilon}}\rangle \simeq (1+m)\langle\bm{g}\rangle + \mathbf{c},
\end{equation}
where $m$ is the multiplicative bias and $\mathbf{c}$ is the additive bias. There is a wealth of literature on various methods for measuring and calibrating shear, which stems from this general idea \citep{Hirata2003, Metacal, Metadet, Anacal}. However, there is a fundamental limit for the precision of traditional WL shear measurements due to the degeneracy between shape and shear, called shape noise. The value of shape noise is typically $\sigma_\epsilon^\mathrm{WL} \approx 0.3$ per galaxy \citep{Bernstein2002}, much larger than typical shear signals $g \sim 0.01$.

KL breaks the shape-shear degeneracy by including the galaxy's kinematic information. We only summarize the key concepts here, following the formalism of \citep{KLRoman, KLDESI}; for a more detailed explanation refer to \citep{KL, KLInferenceI, KLRoman}. The spectra of a galaxy most directly constrain its line-of-sight (LoS) velocity field, which is distorted by cosmic shear in the same way as its image:
\begin{equation}
\label{eqn:shear-vfield}
    v_\mathrm{LoS}^\mathrm{obs}(\bm{\theta}) = v_\mathrm{LoS}(\textbf{A}\cdot\bm{\theta}).
\end{equation}
Here, $v_\mathrm{LoS}$ is the velocity field before lensing, and $\textbf{A}$ is the shear distortion matrix that maps observed positions to positions in the source plane,
\begin{equation}
\label{eqn:shear-mat}
     \textbf{A}=
    (1-\kappa)\begin{pmatrix}
    1-g_1 & -g_2 \\
    -g_2 & 1+g_1
    \end{pmatrix}.
\end{equation}

It is convenient to express shear in the frame defined by the galaxy's major and minor axes, $(g_1, g_2) \rightarrow (g_+, g_\times)$, where $g_+$ stretches the galaxy along its major or minor axis and $g_\times$ stretches it at $45\degree$ to its axes. The transformation is
\begin{equation}
\label{eqn:shear-coords}
    \begin{aligned}
    g_+ = g_1\cos(2\theta_{int}) + g_2\sin(2\theta_{int}), \\
    g_\times = -g_1\sin(2\theta_{int}) + g_2\cos(2\theta_{int}),
    \end{aligned}
\end{equation}
where $\theta_{int}$ is the intrinsic position angle of the galaxy, measured counterclockwise from the $x$-axis. To first order in shear, the maximum LoS velocities of a disk galaxy along its observed major and minor axes are then
\begin{equation}
\label{eqn:v-measured}
    \begin{aligned}
    v^\prime_\mathrm{major}&=-\vcirc\,\sin{i},\\
    v^\prime_\mathrm{minor}&=
    \vcirc\,\sin{i}\,\cos{i}\,\left(1+\frac{1+e_\mathrm{obs}^2}{2e_{\mathrm{int}}}\right)g_\times\,,
    \end{aligned}
\end{equation}
where $i$ is the inclination angle of the galaxy, $\vcirc$ is its maximum circular velocity, and $e\equiv|\hat{\bm{\epsilon}}|$. Without shear, the minor axis carries no LoS velocity, so $v^\prime_\mathrm{minor}$ directly traces $g_\times$. The intrinsic ellipticity follows from the inclination through the geometry of the disk; for a disk whose thickness is a fraction $q_z$ of its radius,
\begin{equation}
\label{eqn:eint-geo}
    e_{\mrm{int}}=\frac{1-\sqrt{1-(1-q_z^2)\sin{^2i}}}{1+\sqrt{1-(1-q_z^2)\sin{^2i}}}\,.
\end{equation}
We adopt $q_z=0.1$ throughout this work, as in \citep{KLDESI}. Finally, $\vcirc$ is tied to the galaxy's luminosity through the TFR \citep{TFR},
\begin{equation}
\label{eqn:TFR}
    \log_\mathrm{10}(\vcirc)=a + bM+\epsilon_\mathrm{TF}\,,
\end{equation}
with intercept $a$, slope $b$, and absolute magnitude $M$. The term $\epsilon_\mathrm{TF}$ is the intrinsic scatter about the relation, which we take to be Gaussian with standard deviation $\sigma_\mathrm{TF}$.

Combining Eqs.~(\ref{eqn:shear-vfield}--\ref{eqn:TFR}) gives us the following estimators for shear \citep{KL, KLDESI}:
\begin{equation}
\label{eqn:reduced_shears}
    \begin{aligned}
    \hat{g}_+ &= \frac{e_\mrm{obs}^2-e_\mrm{int}^2}{2e_\mrm{obs}^2(1-e_\mrm{int}^2)},\\
    \hat{g}_\times&= \left|\frac{v_\mrm{minor}^\prime}{v_\mrm{major}^\prime}\right|\frac{2e_\mrm{int}}{\cos{i}\,(2e_\mrm{int}+1+e_\mrm{obs}^2)}\,.
    \end{aligned}
\end{equation}
We can measure $e_\mathrm{obs}$ from the image and $v^\prime_\mathrm{major}$, $v^\prime_\mathrm{minor}$ from the spectra. The inclination, and therefore $e_\mathrm{int}$, is constrained by comparing $v^\prime_\mathrm{major}$ with the $\vcirc$ expected from the TFR. In practice, past works have performed KL measurements by jointly fitting the relevant galaxy morphology and kinematic parameters with a forward model, rather than using Eq.~(\ref{eqn:reduced_shears}) directly \citep{KLInferenceI, KLRoman, KLDESI}. In this work, we instead estimate the posterior of these parameters with a neural network trained on simulated images and spectra.
